\clearpage
\section{Wstęp}
\label{sec:wstep}

\subsection{Wprowadzenie do problematyki}
\label{subsec:wprowadzenie}

Zatory drogowe wydłużają czas przejazdu, a związany z nimi ruch z częstym zatrzymywaniem i ruszaniem zwiększa zużycie paliwa oraz emisję zanieczyszczeń \cite{treiber2008congestion}.

Jednym z podstawowych narzędzi zarządzania ruchem w miastach jest sterowanie sygnalizacją świetlną na skrzyżowaniach. Sposób przydzielania prawa przejazdu poszczególnym strumieniom pojazdów bezpośrednio wpływa na przepustowość skrzyżowania, czas oczekiwania oraz długość kolejek, a tym samym na sprawność całej sieci drogowej \cite{LEITNER2022507}.

Tradycyjne programy sygnalizacji często działają według z góry ustalonych reguł. W sterowaniu stałoczasowym sekwencję faz i czas ich trwania określa się z góry, bez uwzględniania bieżącej sytuacji ruchowej. Bardziej zaawansowane metody, takie jak sterowanie akomodacyjne typu gap-based, reagują na obecność i napływ pojazdów, ale nadal opierają się na ustalonej logice, a nie na polityce wyuczonej na podstawie doświadczenia.

Alternatywą jest adaptacyjne sterowanie wykorzystujące uczenie ze wzmocnieniem. Agent uczy się polityki przez interakcję ze środowiskiem: obserwuje stan skrzyżowania, wybiera fazę sygnalizacji i otrzymuje nagrodę oceniającą skutki tej decyzji. Polityka nie musi więc być z góry zaprogramowana, ponieważ powstaje podczas uczenia. Potencjał tego podejścia jest punktem wyjścia pracy.

\subsection{Problem badawczy i hipoteza}
\label{subsec:problem-hipoteza}

Sformułowano następujące pytanie badawcze:
\begin{quote}
    \emph{Czy algorytm głębokiego uczenia ze wzmocnieniem może sterować sygnalizacją świetlną na pojedynczym skrzyżowaniu skuteczniej niż sterowanie stałoczasowe oraz akomodacyjne?}
\end{quote}

Przyjęto również następującą hipotezę:
\begin{quote}
    \emph{Odpowiednio wytrenowany algorytm DQN może zmniejszyć opóźnienie, czas oczekiwania oraz długość kolejek pojazdów względem klasycznych metod sterowania, szczególnie w warunkach zmiennego lub wysokiego natężenia ruchu.}
\end{quote}

Hipoteza dotyczy wskaźników opisujących płynność ruchu, przede wszystkim opóźnienia, czasu oczekiwania oraz długości kolejek. Ich poprawa nie musi jednak oznaczać przewagi DQN według każdej metryki ani w każdym scenariuszu. Może się na przykład wiązać z częstszym przełączaniem sygnalizacji. Zakres i warunki potwierdzenia hipotezy oceniono w dalszej części pracy, porównując metody sterowania na odrębnym zbiorze testowym.

\subsection{Cel i zakres pracy}
\label{subsec:cel-zakres}

Celem pracy jest zaprojektowanie, wytrenowanie oraz ocena algorytmu DQN sterującego sygnalizacją świetlną na zamodelowanym skrzyżowaniu ulicy Belgradzkiej i alei Komisji Edukacji Narodowej (al. KEN).

Aby osiągnąć ten cel, zrealizowano następujące cele szczegółowe:
\begin{itemize}
    \item odwzorowanie skrzyżowania w środowisku symulacyjnym SUMO;
    \item implementacja środowiska uczenia ze wzmocnieniem, integrującego symulację z procesem treningowym;
    \item zdefiniowanie przestrzeni obserwacji, przestrzeni akcji oraz funkcji nagrody;
    \item przygotowanie rozdzielnych zbiorów: treningowego, walidacyjnego oraz testowego;
    \item implementacja metod referencyjnych: sterowania stałoczasowego oraz akomodacyjnego;
    \item wyznaczenie odpowiedniej liczby kroków treningowych;
    \item porównanie dwóch sposobów generowania danych treningowych;
    \item dobór hiperparametrów algorytmu DQN;
    \item ocena stabilności procesu uczenia;
    \item porównanie rozpatrywanych metod na odrębnym zbiorze testowym.
\end{itemize}

\subsection{Założenia i ograniczenia pracy}
\label{subsec:zalozenia-ograniczenia}

Zakres pracy świadomie ograniczono, aby przeprowadzić kontrolowany eksperyment. Badano pojedyncze skrzyżowanie bez uwzględniania wpływu skrzyżowań sąsiednich. W modelu ruchu uwzględniono wyłącznie pojazdy silnikowe, pomijając ruch pieszy, rowerowy i komunikację miejską. Ruch drogowy wygenerowano syntetycznie za pomocą skryptu \texttt{randomTrips.py} \cite{sumo_randomtrips_docs}, dlatego nie odwzorowuje on bezpośrednio rzeczywistych pomiarów natężenia. Wszystkie eksperymenty przeprowadzono wyłącznie w symulacji, bez wdrożenia na rzeczywistym sterowniku sygnalizacji.

Celem pracy nie było stworzenie gotowego systemu produkcyjnego ani certyfikowanego sterownika ruchu. Przeprowadzono kontrolowany eksperyment, aby ocenić potencjał DQN w adaptacyjnym sterowaniu sygnalizacją świetlną w porównaniu z klasycznymi, deterministycznymi metodami sterowania.
